\def\pgfsysdriver{pgfsys-dvipdfm.def} % required for tress inside
\documentclass[aspectratio=169
  , xcolor={svgnames}
  , hyperref={ colorlinks,citecolor=Blue
             , linkcolor=DarkRed,urlcolor=DarkBlue}
  , russian
  ]{beamer}
\usetheme{CambridgeUS}
\usefonttheme{professionalfonts}

\usepackage{pgfpages}

\usepackage{bibentry}
\usepackage{cite}
\def\newblock{\hskip .11em plus .33em minus .07em}


\makeatletter
\@ifclassloaded{beamer}{
  % get rid of header navigation bar
  \setbeamertemplate{headline}{}
  % get rid of bottom navigation symbols
  \setbeamertemplate{navigation symbols}{}
  % get rid of footer
  %\setbeamertemplate{footline}{}
}
{}
\makeatother
%%%%%%%%%%%%%%%%%%%%%%%%%%%%%%%%%%%%%%%%%%%%%
\usepackage{fontawesome}
% \newfontfamily{\FA}{Font Awesome 5 Free} % some glyphs missing
\expandafter\def\csname faicon@facebook\endcsname{{\FA\symbol{"F09A}}}
\def\faQuestionSign{{\FA\symbol{"F059}}}
\def\faQuestion{{\FA\symbol{"F128}}}
\def\faExclamation{{\FA\symbol{"F12A}}}
\def\faUploadAlt{{\FA\symbol{"F093}}}
\def\faLemon{{\FA\symbol{"F094}}}
\def\faPhone{{\FA\symbol{"F095}}}
\def\faCheckEmpty{{\FA\symbol{"F096}}}
\def\faBookmarkEmpty{{\FA\symbol{"F097}}}

\newcommand{\faGood}{\textcolor{ForestGreen}{\faThumbsUp}}
\newcommand{\faBad}{\textcolor{red}{\faThumbsODown}}
\newcommand{\faWrong}{\textcolor{red}{\faTimes}}
\newcommand{\faMaybe}{\textcolor{blue}{\faQuestion}}
\newcommand{\faCheckGreen}{\textcolor{ForestGreen}{\faCheck}}
%%%%%%%%%%%%%%%%%%%%%%%%%%%%%%%%%%%%%%%%%%%%%

\usepackage{fontspec}
\usepackage{xunicode}
\usepackage{xltxtra}
\usepackage{xecyr}
\usepackage{hyperref}

\setmainfont[
 Ligatures=TeX,
 Extension=.otf,
 BoldFont=cmunbx,
 ItalicFont=cmunti,
 BoldItalicFont=cmunbi,
% Scale = 1.1
]{cmunrm}
\setsansfont[
 Ligatures=TeX,
 Extension=.otf,
 BoldFont=cmunsx,
 ItalicFont=cmunsi,
%  Scale = 1.2
]{cmunss}
%\setmainfont[Mapping=tex-text]{DejaVu Serif}
%\setsansfont[Mapping=tex-text]{DejaVu Sans}
%\setmonofont{Fira Code}[Contextuals=AlternateOff]
\setmonofont{Fira Code}[Contextuals=Alternate,Scale=0.9]
\newfontfamily{\myfiracode}[Scale=1.5,Contextuals=Alternate]{Fira Code}
%\setmonofont[Scale=0.9,BoldFont={Inconsolata Bold}]{Inconsolata}

\usepackage{polyglossia}
\setmainlanguage{russian}
\setotherlanguage{english}


%\newfontfamily\dejaVuSansMono{DejaVu Sans Mono}
% https://github.com/vjpr/monaco-bold/raw/master/MonacoB/MonacoB.otf
%\newfontfamily\monacoB{MonacoB}
%%%%%%%%%%%%%%%%%%%%%%%%%%%%%%%%%%%%%%%%%%%%%%%5
\usepackage{soul} % for \st that strikes through
\usepackage[normalem]{ulem} % \sout

\usepackage{stmaryrd}
\newcommand{\sem}[1]{\ensuremath{\llbracket #1\rrbracket}}


\usepackage{listings}
%\lstdefinestyle{style1}{
%  language=haskell,
%  numbers=left,
%  stepnumber=1,
%  numbersep=10pt,
%  tabsize=4,
%  showspaces=false,
%  showstringspaces=false
%}
%\lstdefinestyle{hsstyle1}
%{ language=haskell
%%          , basicstyle=\monacoB
%         , deletekeywords={Int,Float,String,List,Void}
%         , breaklines=true
%         , columns=fullflexible
%         , commentstyle=\color{ForestGreen}
%         , escapeinside=§§
%         , escapebegin=\begin{russian}\commentfont
%         , escapeend=\end{russian}
%         , commentstyle=\color{ForestGreen}
%         , escapeinside=§§
%         , escapebegin=\begin{russian}\color{ForestGreen}
%         , escapeend=\end{russian}
%         , mathescape=true
%%          , backgroundcolor = \color{MyBackground}
%}
%
%\newcommand{\inline}[1]{\lstinline{haskell}{#1}}
%\def\hsinline{\mintinline{haskell}}
%\def\inline{\hsinline}
%
%\lstnewenvironment{hslisting} {
%    \lstset { style={hsstyle1} }
%  }
%  {}
%  
%%%%%%%%%%%%%%%%%%%%%%%%%%%%%%%%%%%%%%%%%%%%%%%%%%%%%%%%%%%  
%%\setmainfont[
%% Ligatures=TeX,
%% Extension=.otf,
%% BoldFont=cmunbx,
%% ItalicFont=cmunti,
%% BoldItalicFont=cmunbi,
%%]{cmunrm}
%%% С засечками (для заголовков)
%%\setsansfont[
%% Ligatures=TeX,
%% Extension=.otf,
%% BoldFont=cmunsx,
%% ItalicFont=cmunsi,
%%]{cmunss}
%% \setmonofont[Scale=0.6]{Monaco}
%
%\usefonttheme{professionalfonts}
%\usepackage{times}
\usepackage{tikz}
\usetikzlibrary{cd}
\usepackage{tikz-cd}
\usepackage{caption}
\usepackage{subcaption}

%\renewtheorem{definition}{برهان}[chapter]
%%\DeclareMathOperator{->}{\rightarrow}
%\newcommand\iso{\ensuremath{\cong}}
%\usepackage{verbatim}
%\usepackage{graphicx}
%\usetikzlibrary{arrows,shapes}

%\usepackage{amsmath}
%\usepackage{amsfonts}
\usepackage{scalerel}
\DeclareMathOperator*{\myvee}{\scalerel*{\vee}{\sum}}
\DeclareMathOperator*{\mywedge}{\scalerel*{\wedge}{\sum}}

%
%\usepackage{tabulary}
%
%% sudo aptget install ttf-mscorefonts-installer
%%\setmainfont{Times New Roman}
%%\setsansfont[Mapping=tex-text]{DejaVu Sans}
%
%%\setmonofont[Scale=1.0,
%%    BoldFont=lmmonolt10-bold.otf,
%%    ItalicFont=lmmono10-italic.otf,
%%    BoldItalicFont=lmmonoproplt10-boldoblique.otf
%%]{lmmono9-regular.otf}
%
\usepackage[cache=true]{minted}
\usemintedstyle{perldoc}

\def\hsinline{\mintinline{haskell}}
\def\mlinline{\mintinline{ocaml}}
% color options
\definecolor{YellowGreen} {HTML}{B5C28C}
\definecolor{ForestGreen} {HTML}{009B55}
\definecolor{MyBackground}{HTML}{F0EDAA}



\institute{матмех СПбГУ}

\addtobeamertemplate{title page}{}{
  \begin{center}{\tiny Дата сборки: \today}\end{center}
}


\usepackage{subcaption}
\usepackage{etoolbox}


\usepackage{exercise}
\usepackage{tikz}
 \usepackage{tikz-qtree}
\usetikzlibrary{trees}
\usepackage[edges]{forest}
\forestset{.style={
%  for tree={l=1em, l sep=1em, s sep=1em}
  forked edges,
    for tree={    grow'=0,    draw,    align=c,    font=\sffamily,
        rounded corners  }
  }}

\newcommand{\lstquot}[1]{``\lstinline{#1}''}
\newcommand{\sembr}[1]{\llbracket{#1}\rrbracket}
\newcommand\false{$f\!alse$}
\newcommand\myif{i\!f}


\def\transarrow{\xrightarrow}
\newcommand{\setarrow}[1]{\def\transarrow{#1}}

\def\padding{\phantom{X}}
\newcommand{\setpadding}[1]{\def\padding{#1}}

\def\subarrow{}
\newcommand{\setsubarrow}[1]{\def\subarrow{#1}}

\newcommand{\trule}[2]{\dfrac{#1}{#2}}
\newcommand{\crule}[3]{\dfrac{#1}{#2},\;{#3}}
\newcommand{\withenv}[2]{{#1}\vdash{#2}}
\newcommand{\trans}[3]{{#1}\transarrow{\padding{\textstyle #2}\padding}\subarrow{#3}}
\newcommand{\ctrans}[4]{{#1}\transarrow{\padding#2\padding}\subarrow{#3},\;{#4}}
\newcommand{\llang}[1]{\mbox{\lstinline[mathescape]|#1|}}
\newcommand{\pair}[2]{\inbr{{#1}\mid{#2}}}
\newcommand{\inbr}[1]{\left<{#1}\right>}
\newcommand{\highlight}[1]{\color{red}{#1}}
%\newcommand{\ruleno}[1]{\eqno[\scriptsize\textsc{#1}]}
\newcommand{\ruleno}[1]{\mbox{[\textsc{#1}]}}
\newcommand{\rulename}[1]{\textsc{#1}}
\newcommand{\inmath}[1]{\mbox{$#1$}}
%\newcommand{\lfp}[1]{fix_{#1}}
%\newcommand{\gfp}[1]{Fix_{#1}}
\newcommand{\vsep}{\vspace{-2mm}}
\newcommand{\supp}[1]{\scriptsize{#1}}
\renewcommand{\sembr}[1]{\llbracket{#1}\rrbracket}
\newcommand{\cd}[1]{\texttt{#1}}
%\newcommand{\free}[1]{\boxed{#1}}
%\newcommand{\binds}{\;\mapsto\;}
%\newcommand{\dbi}[1]{\mbox{\bf{#1}}}
%\newcommand{\sv}[1]{\mbox{\textbf{#1}}}
%\newcommand{\bnd}[2]{{#1}\mkern-9mu\binds\mkern-9mu{#2}}
%\newcommand{\meta}[1]{{\mathcal{#1}}}
%\newcommand{\dom}[1]{\mathtt{dom}\;{#1}}
%\newcommand{\primi}[2]{\mathbf{#1}\;{#2}}
%\renewcommand{\dom}[1]{\mathcal{D}om\,({#1})}
%\newcommand{\ran}[1]{\mathcal{VR}an\,({#1})}
%\newcommand{\fv}[1]{\mathcal{FV}\,({#1})}
%\newcommand{\tr}[1]{\mathcal{T}r_{#1}}
\newcommand{\diseq}{\not\equiv}
%\newcommand{\reprfunset}{\mathcal{R}}
%\newcommand{\reprfun}{\mathfrak{f}}
%\newcommand{\cstore}{\Omega}
%\newcommand{\cstoreinit}{\cstore_\epsilon^{init}}
%\newcommand{\csadd}[3]{add(#1, #2 \diseq #3)}  %{#1 + [#2 \diseq #3]}
%\newcommand{\csupdate}[2]{update(#1, #2)}  %{#1 \cdot #2}
\newcommand{\primi}[1]{\mathbf{#1}}
\newcommand{\sem}[1]{\llbracket #1 \rrbracket}
\newcommand{\ir}{\ensuremath{\mathcal{S}}}
\usepackage{tikz}
\newcommand*\circled[1]{\tikz[baseline=(char.base)]{
    \node[shape=circle,draw,inner sep=1pt] (char) {#1};}}

\let\emptyset\varnothing
\let\eps\varepsilon

% for fancy table
\newcommand{\lheadl}[2]{\multicolumn{1}{|>{\centering\arraybackslash}m{#1}|}{{#2}}}
\newcommand{\head}[2]{\multicolumn{1}{>{\centering\arraybackslash}m{#1}}{\textbf{\small #2}}}
\newcommand{\headll}[2]{\multicolumn{1}{>{\centering\arraybackslash}m{#1}||}{\textbf{\small #2}}}
\newcommand{\lheadll}[2]{\multicolumn{1}{|>{\centering\arraybackslash}m{#1}||}{\textbf{\small #2}}}
\newcommand{\headl}[2]{\multicolumn{1}{>{\centering\arraybackslash}m{#1}|}{\textbf{\small #2}}}
\usepackage{longtable}
\newcommand{\nodata}{}
\newcommand{\tablenotemark}[1]{#1}

\newcommand{\contributions}{
\begin{enumerate}
\item[I] Synthesis via combination of two relational interpreters
\item[II] Replacing $\forall$ quantifier by a finite set of examples
\item[III] Branch and bound optimization of synthesis using \emph{structural constraints}
%\item[IV] Extension of OCanren called \emph{structural constraint}
\end{enumerate}
}

%\newcommand{\treeboolsexample}[1] {
%\begin{lstlisting}[basicstyle=\small,belowskip=1em,aboveskip=0em]
%(_, false, true) -> 1
%(false, true, _) -> 2
%(_, _, false) -> 3
%(_, _, true) -> 4
%\end{lstlisting}
%}
%%%%%%%%%%%%%%%%%%%%%%%%%%%%%%%%%%%%%%%%



\title{Реляционный синтез сопоставления с образцом}
\subtitle{Relational Synthesis for Pattern Matching}

\date{9 ноября 2020}
\author{Косарев Дмитрий} 
\institute[]{\normalfont
Будет опубликовано на \\
Asian Symposium on Programming Languages and Systems (APLAS) 2020
%\textsuperscript{1} JetBrains Research, Programming Languages and Tools Lab
% \and \inst{2} Saint-Petersburg University
 }


\newcommand{\verbatimfont}[1]{\def\verbatim@font{#1}}
\usepackage{verbatimbox}

\forestset{
    .style={
        for tree={
            base=bottom,
            child anchor=north,
            align=center,
            s sep+=1cm,
    straight edge/.style={
        edge path={\noexpand\path[\forestoption{edge},thick,-{Latex}] 
        (!u.parent anchor) -- (.child anchor);}
    },
    if n children={0}
        {tier=word, draw, thick, rectangle}
        {draw, diamond, thick, aspect=2},
    if n=1{%
        edge path={\noexpand\path[\forestoption{edge},thick,-{Latex}] 
        (!u.parent anchor) -| (.child anchor) node[pos=.2, above] {Y};}
        }{
        edge path={\noexpand\path[\forestoption{edge},thick,-{Latex}] 
        (!u.parent anchor) -| (.child anchor) node[pos=.2, above] {N};}
        }
        }
    }
}
  

\begin{document}
{
\begin{frame}[fragile]
  \begin{tabular}{p{5.5cm} p{5.5cm}}
   \begin{center}
%      \includegraphics[height=1.5cm]{pictures/jetbrainsResearch.pdf}
    \end{center}
    &
    \begin{center}
   %   \includegraphics[height=1.5cm]{pictures/SPbGU_Logo.png}
    \end{center}
  \end{tabular}
  \titlepage
\end{frame}
}

%\maketitle

% For every picture that defines or uses external nodes, you'll have to
% apply the 'remember picture' style. To avoid some typing, we'll apply
% the style to all pictures.
\tikzstyle{every picture}+=[remember picture] 

% By default all math in TikZ nodes are set in inline mode. Change this to
% displaystyle so that we don't get small fractions.
\everymath{\displaystyle}
%\begin{comment}
% Uncomment these lines for an automatically generated outline.
\begin{frame}{Pattern Matching}
Efficient compilation of pattern matching is essential for functional programming
\vspace{1cm}

%\includegraphics{pictures/SPbGU_Logo.png}

Два основных подхода:
\begin{itemize}
\item диаграммы решений
\begin{itemize}
\item минимизируется количество проверок
\end{itemize}
\item автомат с возвратами
\begin{itemize}
\item минимизируется размер кода
\end{itemize}%\pause
\item \textbf{synthesis} (our approach)
\end{itemize}
% \note{}

\end{frame}


\defverbatim[colored]{\matchA}{
\begin{lstlisting}
match x,y,z with
| _,F,T -> 1
| F,T,_ -> 2
| _,_,F -> 3
| _,_,T -> 4
\end{lstlisting}
}

\defverbatim[colored]{\matchB}{
\begin{lstlisting}
if x then
  if y then
    if z then 4 else 3
  else
    if z then 1 else 3
else
  if y then 2
  else
    if z then 1 else 3
\end{lstlisting}
}

\defverbatim[colored]{\matchC}{
\begin{lstlisting}
if y then
  if x then
    if z then 4 else 3
  else 2
else
  if z then 1 else 3
\end{lstlisting}
}


\defverbatim[colored]{\forestB}{
\begin{forest} 
[$x$%, tikz={\draw[{Latex}-, thick] (.north) --++ (0,1);}
    [$y$
          [$z$
                [4] 
                [3] 
            ]   
         [$z$
               [1] 
               [3] 
           ]    
    ]   
    [$y$
        [2] 
        [$z$
            [1] 
            [3] 
        ]   
    ]   
] 
\end{forest}
}



\begin{frame}[fragile]{An Example of Pattern Matching and It's Two Compiled Representations}

\begin{minipage}[t]{0.25\linewidth}
\begin{minipage}{7cm}
\matchA
\end{minipage}
%\caption{Pattern matching}
\end{minipage}
\hspace{0.5cm}
\begin{minipage}[t]{0.32\linewidth}
\only<1>{\begin{minipage}{7cm}
\matchB
\end{minipage}}
\only<2>{\begin{minipage}{7cm}
\matchC
\end{minipage}}

\end{minipage}
\hspace{0.5cm}
\begin{minipage}[t]{0.3\linewidth}
\only<1>{\begin{minipage}{7cm}
\forestB
\end{minipage}}
\only<2>{\begin{minipage}{7cm}
asdf
\end{minipage}}
\end{minipage}
%\caption{Pattern matching compilation can be non-trivial (example from~\cite{maranget2008}).}\label{fig:match-example}
\end{frame}

\begin{frame}[fragile]
\begin{forest} 
[$x$%, tikz={\draw[{Latex}-, thick] (.north) --++ (0,1);}
    [$y$
          [$z$
                [4] 
                [3] 
            ]   
         [$z$
               [1] 
               [3] 
           ]    
    ]   
    [$y$
        [2] 
        [$z$
            [1] 
            [3] 
        ]   
    ]   
] 
\end{forest}
\end{frame}

\begin{frame}[fragile]{Пример: Деревья и DAGи терма $f(g(a, x), g(a,x))$}
\begin{center}
\begin{tikzpicture}[->,thick,scale=0.5, every node/.style={scale=0.5}
  , grow via three points={
      one child at (0,-2) and 
      two children at (.8,-2) and (-0.8,-2) 
  }]
     \tikzstyle{tnode}=[circle, inner sep=1.5mm]
     \tikzstyle{var}=[ circle, inner sep=1.5mm,draw]
     \def\rstep{5cm}
     \huge
         
    \begin{scope}[xshift= \rstep]
      \node[tnode] {f}
          child {node[tnode] {g} 
            child {node[var] {x} }
            child {node[tnode] {a} }
          }
          child[missing] {} 
          child {node[tnode] {g} 
            child {node[var] {x} }
            child {node[tnode] {a} }
          };
%        child[missing] {} ;
      \end{scope}
      
    \begin{scope}[xshift=15cm]
      \node[tnode] {f}
          child {node[tnode] (gggg) {g} 
            child[missing] {node[var] {x} }
            child {node[tnode]  {a} }
          }
          child[missing] {}
          child {node[tnode]  {g} 
            child {node[var]  (xxxx) {x} }
            child {node[tnode] {a} }
          };

      \draw [->,thick](gggg.south)
              to [out=-90,in=90] (2,-5)
              to [out=-90,in=0] (1,-6) 
              to [out=180,in=-45]  (xxxx.south); 
      \end{scope}
      
      \begin{scope}[xshift=23cm]
        \node[tnode] (ffff) {f}
            child {node[tnode] (gggg) {g} 
              child {node[var] {x} }
              child {node[tnode] {a} }
            };

        \draw [->,thick](ffff)
                to [out=0,in=90] (1,-1)
                to [out=-90,in=45]  (gggg.east); 
        \end{scope}
 \end{tikzpicture}
\end{center}
\end{frame}

\begin{frame}{Contributions}
\Large
\contributions
\end{frame}


\begin{frame}[fragile]{Pattern Matching Syntax}
\begin{figure}[ht]
\begin{subfigure}[t]{0.4\linewidth}
$$
 \begin{array}{rcll}
    \mathcal{C} & = & \{ C_1^{k_1}, \dots, C_n^{k_n} \}\\
    \mathcal{V} & = & \mathcal{C}\,\mathcal{V}^*\\  
    \mathcal{P} & = & \_ \mid \mathcal{C}\,\mathcal{P}^*
 \end{array}
$$
%TODO: paterrn matching task
$$
 \begin{array}{c}
\trans{\inbr{v;\,p_1,\dots,p_k}}{}{i}\\
1\leqslant i\leqslant k+1
 \end{array}
$$
\end{subfigure}
\hspace{0.5cm}
\begin{subfigure}[t]{0.5\linewidth}
\[
\begin{array}{rcl}
  \mathcal M & = & \bullet \\
  &   & \mathcal M\,[\mathbb{N}] \\
  \ir & = & \primi{return}\,\mathbb{N} \\
  &   & \primi{switch}\;\mathcal{M}\;\primi{with}\; [\mathcal{C}\; \primi{\rightarrow}\; \ir]^*\;\primi{otherwise}\;\ir
\end{array}
\]
\end{subfigure}
\end{figure}
\vspace{0.5cm}
Things that are omitted for simplicity: guards, pattern variables

%TODO: program size here?
\end{frame}

\begin{frame}{Synthesis Algorithm}
%High-level description:
\[
%depth^o\,v\,n \;\; \wedge \;\;
\forall v \quad \forall  (1\leqslant\!i\leqslant \!k\!+\!1) \quad
%\mbox{\lstinline|fresh ($i$)|}\; \{
(match^o\,v\,\,p_1,\dots,p_k\,\,i) \;\; \Leftrightarrow \;
eval^o_{\mathcal S}\,v\,\circled{?}\,i
%\}
\]

\begin{itemize}
%\item $depth^o\,v\,n$ generates examples bound by depth $n$
%\item $\forall v$ -- all possible scrutinees
%\item $\forall i$ -- all possible branches
\item $match^o\,v\,\,p_1,\dots,p_k\,\,i$ -- a declarative interpreter of pattern matching, for every  possible scrutinee $v$ returns a single branch $i$ 
\item $\circled{?}$ is a program to be synthesized
\item $eval^o_{\mathcal S}\,v\,\circled{?}\,i$ -- a declarative interpreter of $\mathcal S$, ensures that synthesized program $\circled{?}$ on example $v$ returns a correct branch index $i$
\end{itemize}
\vspace{1cm}
\begin{itemize}
\item[\faGood] It's easy to implement $match^o$ and  $eval^o$ (even for extension of pattern matching)
\item[\faBad]  No universal quantifier in miniKanren with disequality constraints
\end{itemize}

\end{frame}


\begin{frame}{2Eliminating $\forall$ by Generating a Finite Set of Examples}
For every concrete pattern matching we know:
\begin{itemize}
\item the type of scrutinee
\item all its patterns 
\begin{itemize}
\item maximal depth of these patterns
\end{itemize}
\end{itemize}

\begin{alertblock}{\textbf{Idea}}
Let's generate all inhabitants of the scrutinee's type until certain depth and use them as examples
\end{alertblock}
\vspace{1cm}

\begin{itemize}
%\item[\faGood] Is correct (proved in Coq)
\item[\faBad] The number of examples is  worst case exponential
\end{itemize}
\end{frame}

\begin{comment}
\begin{frame}{Our Method: Synthesis + Relational Interpreter}
Do \emph{not compile} with specific algorithm but \emph{synthesize} compiled representation on large enough but finite set of examples
\vspace{1cm}

Fast $\Leftrightarrow$ Small: synthesized program should contain less checks.
\vspace{1cm}

We are using relational programming, more precisely OCanren~\cite{OCanrenWeb} from miniKanren~\cite{MiniKanrenWeb} family.
\vspace{1cm}


Our repo on Github: ~\cite{Repo}.
\end{frame}
\end{comment}



\begin{frame}[fragile]{Example 1: Three Examples is a Complete Set}
\centering
\begin{minipage}{.381\textwidth}
\begin{lstlisting}{ocaml}
match (s : unit list) with 
| _ :: _ :: _ -> 1
| _           -> 2
\end{lstlisting}
\end{minipage}

\begin{figure}
\begin{subfigure}[t]{0.3\linewidth}
\begin{subfigure}[b][4cm][t]{0.2\linewidth}
\begin{tikzpicture}[sibling distance=3em,
  every node/.style = {shape=rectangle, rounded corners, minimum height=.5cm,
    draw, align=center, font=\ttfamily
    ,top color=white, bottom color=blue!20}]]
  \node {Cons}
    child { node {\phantom{1}\_\phantom{1}} }
    child { node {Cons}
      child { node  {\phantom{1}\_\phantom{1}} }
      child { node {\phantom{1}\_\phantom{1}} } 
    };
\end{tikzpicture}
\end{subfigure}
\hspace{2cm}
\begin{subfigure}[b][4cm][t]{0.1\linewidth}
\begin{tikzpicture}[sibling distance=3em,
  every node/.style = {shape=rectangle, rounded corners, minimum height=.5cm,
    draw, align=center, font=\ttfamily
    ,top color=white, bottom color=blue!20}]]
  \node {\phantom{1}\_\phantom{1}};
\end{tikzpicture}
\end{subfigure}
\caption{Patterns}
\end{subfigure}
\hspace{1.5cm}
\begin{subfigure}[t]{0.5\linewidth}
\begin{subfigure}[b][4cm][t]{0.2\linewidth}
\begin{tikzpicture}[sibling distance=3em,
  every node/.style = {shape=rectangle, rounded corners, minimum height=.5cm,
    draw, align=center, font=\ttfamily
    ,top color=white, bottom color=green!20}]]
  \node {Cons}
    child { node {()} }
    child { node {Cons}
      child { node {()} }
      child { node { Nil } } 
    };
\end{tikzpicture}
\end{subfigure}
\hspace{1cm}
\begin{subfigure}[b][4cm][t]{0.15\linewidth}
\begin{tikzpicture}[sibling distance=3em,
  every node/.style = {shape=rectangle, rounded corners, minimum height=.5cm,
    draw, align=center, font=\ttfamily
    ,top color=white, bottom color=green!20}]]
  \node {Cons}
    child { node {()} }
    child { node { Nil} } ;
\end{tikzpicture}
\end{subfigure}
\hspace{1cm}
\begin{subfigure}[b][4cm][t]{0.15\linewidth}
\begin{tikzpicture}[sibling distance=3em,
  every node/.style = {shape=rectangle, rounded corners, minimum height=.5cm,
    draw, align=center, font=\ttfamily
    , top color=white, bottom color=green!20}]]
  \node {Nil};
\end{tikzpicture}
\end{subfigure}
\caption{Three examples}
\end{subfigure}
\end{figure}
\end{frame}


\begin{frame}[fragile]{Example 2: Not a Complete Set of Examples}
\begin{minipage}{0.4\linewidth}
\begin{lstlisting}{ocaml}
match (s : unit list) with 
| []  -> 1
| _   -> 2
\end{lstlisting}
\end{minipage}
\begin{minipage}{0.55\linewidth}
%\begin{itemize}
%\item 
Two patterns:\\
\begin{tikzpicture}[sibling distance=3em,
  every node/.style = {shape=rectangle, rounded corners, minimum height=.5cm,
    draw, align=center, font=\ttfamily
    ,top color=white, bottom color=blue!20}]]
  \node { Nil };
\end{tikzpicture}
\begin{tikzpicture}[sibling distance=3em,
  every node/.style = {shape=rectangle, rounded corners, minimum height=.5cm,
    draw, align=center, font=\ttfamily
    ,top color=white, bottom color=blue!20}]]
  \node {\phantom{1}\_\phantom{1}};
\end{tikzpicture}
\vspace{1em}

The set of examples bound by depth 1 is a singleton set\\
\begin{tikzpicture}[sibling distance=2em,
  every node/.style = {shape=rectangle, rounded corners, minimum height=.5cm,
    draw, align=center, font=\ttfamily
    , top color=white, bottom color=green!20}]]
  \node {Nil};
\end{tikzpicture}
\vspace{1em}

This program behave well on the set of examples \\
$$
 \begin{array}{ll}
    \primi{switch} \dots \primi{with}   & \\
    |\ \text{Nil} \rightarrow 1  & \\  
    |\ \primi{otherwise} \rightarrow  1 & 
 \end{array}
$$
%\item 
but is wrong 
\end{minipage}
\end{frame}

\begin{comment}
\begin{frame}[fragile]{Example 2 (1/2) }
\centering
\begin{minipage}{.381\textwidth}
\begin{lstlisting}{ocaml}
match (s : unit list) with 
| []  -> 1
| _   -> 2
\end{lstlisting}
\end{minipage}


\begin{figure}
\begin{subfigure}[b][2cm][b]{0.2\linewidth}
\begin{subfigure}[t][.5cm][t]{0.2\linewidth}
\begin{tikzpicture}[sibling distance=3em,
  every node/.style = {shape=rectangle, rounded corners, minimum height=.5cm,
    draw, align=center, font=\ttfamily
    ,top color=white, bottom color=blue!20}]]
  \node { Nil };
\end{tikzpicture}
\end{subfigure}
\hspace{.5cm}
\begin{subfigure}[t][.5cm][t]{0.1\linewidth}
\begin{tikzpicture}[sibling distance=3em,
  every node/.style = {shape=rectangle, rounded corners, minimum height=.5cm,
    draw, align=center, font=\ttfamily
    ,top color=white, bottom color=blue!20}]]
  \node {\phantom{1}\_\phantom{1}};
\end{tikzpicture}
\end{subfigure}
\caption{Two patterns}
\end{subfigure}
\hspace{1cm}
\begin{subfigure}[b]{0.1\linewidth}
%\begin{subfigure}[b][.5cm][b]{0.15\linewidth}
\begin{tikzpicture}[sibling distance=2em,
  every node/.style = {shape=rectangle, rounded corners, minimum height=.5cm,
    draw, align=center, font=\ttfamily
    , top color=white, bottom color=green!20}]]
  \node {Nil};
\end{tikzpicture}
\caption{Single example}
\end{subfigure}
\hspace{1cm}
\begin{subfigure}[b]{0.3\linewidth}
$$
 \begin{array}{ll}
    \primi{switch} \dots \primi{with}   & \\
    |\ \text{Nil} \rightarrow 1  & \\  
    |\ \primi{otherwise} \rightarrow  1 & 
 \end{array}
$$
\caption{Potential incorrect program}
\end{subfigure}
\end{figure}
\end{frame}

\end{comment}

%\begin{figure}
%\begin{subfigure}[b]{0.3\linewidth}
%\begin{forest}
%  forked edges,
%  for tree={    grow'=0,    draw,    align=c,    font=\sffamily,
%      rounded corners  },
%  highlight/.style={    thick,    font=\sffamily\bfseries  }
%    [{pats}
%      [{Nil}]
%      [{\_}]
%    ]
%\end{forest}
%\end{subfigure}
%\hspace{.5cm}
%\begin{subfigure}[b]{0.6\linewidth}
%%\texttt{hello [] 111}
%\begin{lstlisting}[language=ocaml]
%match (s : unit list) with 
%| []  -> 1
%| _   -> 2
%\end{lstlisting}
%\vspace{1cm}
%\begin{itemize}
%\item All inhabitants bound by depth = 1 -- is a set of single element
%\item It is not enough
%\item We should test at least on an expression that starts from Cons  (\lstinline{unit :: []}, for instance)
%\item but if we don't do that...
%\end{itemize}
%\end{subfigure}
%\end{figure}
%\end{frame}

\begin{comment}

\begin{frame}[fragile]{Example 2 (2/2) }
\begin{figure}
\begin{subfigure}[b]{0.3\linewidth}
\begin{forest}
  forked edges,
  for tree={    grow'=0,    draw,    align=c,    font=\sffamily,
      rounded corners  },
  highlight/.style={    thick,    font=\sffamily\bfseries  }
    [{pats}
      [{Nil}]
      [{\_}]
    ]
\end{forest}
\vspace{1cm}

Wrong set of examples:
\begin{itemize}
\item \verb=Nil=
\end{itemize}
\vspace{1cm}
\end{subfigure}
\hspace{.5cm}
\begin{subfigure}[b]{0.6\linewidth}
\begin{lstlisting}{ocaml}
match (s : unit list) with 
| []  -> 1
| _   -> 2
\end{lstlisting}
\vspace{1cm}
\begin{itemize}
\item ... and test single example \verb=[]= (i.e. Nil)
\item then the synthesized program 
\[
 \begin{array}{ll}
    \primi{switch} \dots \primi{with}   & \\
    |\ \text{Nil} \rightarrow 1  & \\  
    |\ \primi{otherwise} \rightarrow  1 & 
 \end{array}
\]
behaves the same on our set of examples
\item but obviously wrong
\end{itemize}
\end{subfigure}
\end{figure}
\end{frame}

\end{comment}

\begin{frame}{Current Algorithm for Examples Generation}
\begin{figure}
\begin{subfigure}[b]{0.75\linewidth}
Algorithm briefly:
\begin{itemize}
\item Evaluate the depth of patterns $h$
\item Synthesize all inhabitants, but
\item on depth $h+1$ use a single predefined per type inhabitant 
\end{itemize}
\vspace{1cm}

\begin{itemize}
%\item[\faGood] Is correct (proved in Coq)
\item[\faBad] The number of examples is  worst case exponential
\end{itemize}
\end{subfigure}
\hspace{.5cm}
\begin{subfigure}[b]{0.45\linewidth}

%\vspace{1cm}
%We are currently working on
%\begin{itemize}
%\item When we can generate less examples?
%\item Overall synthesis speed
%\begin{itemize}
%\item it not only depend on a number of examples
%\end{itemize}
%\end{itemize}
\end{subfigure}
\end{figure}

\end{frame}


\begin{frame}[fragile]{Improvement: Reducing Examples Set}
An example:
\begin{lstlisting}{ocaml}
match (s : bool * bool * bool) with 
| (_,_,F) -> 0
| (_,_,T) -> 1
\end{lstlisting}
\vspace{.5cm}
We can statically determine that 

\begin{itemize}
\item There are $2^3$ inhabitants of type \lstinline=bool * bool * bool=
\item No need to  check that scrutinee is a triple
\item No need to look at the 1st and  the 2nd subtrees of a triple (only the 3rd matters)
\end{itemize}
\vspace{.5cm}


So, we can run the synthesis only on two examples $\{(\mathcal{B},\mathcal{B},\text{T}),(\mathcal{B},\mathcal{B},\text{F})\}$
where $\mathcal{B}$ is any value of type \lstinline=bool=
if we \textbf{forbid the investigation of  some subtrees} of examples


%$\text{T}\!\!\! \text{F}$
\end{frame}


\begin{frame}{Branch and Bound Optimization}
The synthesis answer is stored in variable $\circled{?}$ and 
%Synthesis answer is stored in variable $\circled{?}$ and 
%Our answer is always in single variable -- $\circled{?}$ and 
during the search we \emph{only} extend it with new $\primi{switch}$es \\

\begin{alertblock}{\textbf{Idea}}
If current estimate is longer than the shortest answer already found -- cut this branch
\end{alertblock}

Requires a modified miniKanren's \lstinline=run= primitive: for every answer found we 
\begin{itemize}
\item evaluate it's size
\item update already found minimum
\end{itemize}

Pruning search branches is implemented using a new primitive called  \emph{structural constraint}

\end{frame}

\begin{frame}{Structural constraint}
A new primitive called \emph{structural constraint}:
\begin{itemize}
\item Takes a logic value and performs reification in a current state
\item Takes a predicate which inspects the reified value 
\item If there are too many branches, stops the search ($failure^o$)
\item Otherwise continues search without changing the state ($success^o$)
\end{itemize}
\vspace{1em}
%Peculiarities:
Subtleties:
\begin{itemize}
\item Used to count the size of current estimate
\item May take or may not take into account disequality constraints
\item Can be used to implement $absent^o$ or similar constraints
\item Deals with reified values $\Rightarrow$ can be slow
\end{itemize}
\end{frame}



\begin{frame}[fragile]{Minimization Criterion for Synthesized Program}
\begin{figure}
\begin{subfigure}[b]{0.3\linewidth}
$$
\begin{array}{l}
\primi{switch}\;\mathcal{M}\;\primi{with}\; \\
\mathcal{C}_1\; \primi{\rightarrow}\; \ir_1\\
\dots \\
\mathcal{C}_n\; \primi{\rightarrow}\; \ir_n\\
\primi{otherwise}\;\ir\\
\end{array}\\
$$
$$
\Updownarrow
$$
$$
\begin{array}{l}
\primi{if}\; \mathcal{M} = \mathcal{C}_1\;\primi{then}\; \ir_1\\
\dots \\
\primi{else if}\; \mathcal{M} = \mathcal{C}_n\; \primi{then}\;  \ir_n\\
\primi{else}\;\ir\\
\end{array}\\
$$
\vspace{1cm}

\end{subfigure}
\hspace{1cm}
\begin{subfigure}[b]{0.6\linewidth}
A single switch with $n$ branches roughly can be encoded as $n$ \texttt{if}s\\

Let's say that the size of
\begin{itemize}
\item  $\primi{switch}$ is the number of branches
\item $\primi{return}$ is 0
\item the whole program in $\ir$ is the sum of sizes of its $\primi{switch}$es 
\end{itemize}
Our minimization criterion: reduce the size of synthesized program\\

But there can be other criteria:
\begin{itemize}
\item The  depth 
\item The branching factor
\item etc.
\end{itemize}
\end{subfigure}
\end{figure}

\end{frame}





\begin{comment}


\begin{frame}[fragile,allowframebreaks]{Experiments}
%\begin{table}[H]
  \begin{longtable}{|m{3.5cm}||cccccc|}
\hline
    \lheadll{3cm}{Patterns} & 
%    \head{1.5cm}{Size constraint} & 
%    \headll{1.7cm}{Answers requested} & 
    \head{1.5cm}{\scriptsize \# samples} & 
    \head{1.5cm}{\scriptsize 1st answer size} & 
    \head{1.5cm}{\scriptsize 1st answer time (ms)} & 
    \head{1.5cm}{\scriptsize Answers found} & 
    \head{1.5cm}{\scriptsize Last answer size} & 
    \headl{1.5cm}{\scriptsize Total search time (ms)}
\endhead
    \hline
    \hline

%    \texttt{
%A
%B
%C }&3&&1&2&&2\\ \hline
%
%    \begin{lstlisting}[basicstyle=\scriptsize]
%true
%false
%    \end{lstlisting} &2&&$<$1&1&&$<$1\\
%        \hline

            \begin{lstlisting}[basicstyle=\scriptsize,belowskip=-1.5em,aboveskip=-1em]
(true, _) 
(_, true)
(false, false)
\end{lstlisting}  &4&2&5&1&2&6\\
        \hline
     \begin{lstlisting}[basicstyle=\scriptsize,belowskip=-1.5em,aboveskip=-1em]
(_, false, true)
(false, true, _)
(_, _, false) 
(_, _, true)
\end{lstlisting}   &8&6&$\sim$1000&3&4&$\sim$2000\\
        \hline
%        5&
%     \begin{lstlisting}[basicstyle=\scriptsize]
%([], _)
%(_, [])
%(_ :: _, _ :: _)
%    \end{lstlisting} &100&10&4&5&1&6\\
%        \hline
%        6&
     \begin{lstlisting}[basicstyle=\scriptsize,belowskip=-1.5em,aboveskip=-1em]
type nat = Z 
         | S of nat
  
(Succ _, Succ _)
(Zero, _)
(_, Zero)    \end{lstlisting} &4&2&30&1&2&30
    \\
        \hline
     \begin{lstlisting}[basicstyle=\scriptsize,belowskip=-1.5em,aboveskip=-1em]
type 'a list = 
| Nil 
| Nil2 
| Cons of 'a * 'a list     

(Nil, _)
(_, Nil)
(Nil2, _)
(_, Nil2)
(Cons(_,_),Cons(_,_))
    \end{lstlisting}
     &9&5&24&3& 5  &157                    \\ 
      \hline
%      \hline
     \begin{lstlisting}[basicstyle=\scriptsize,belowskip=-1.5em,aboveskip=-1em]
type code = 
| Push 
|Ldi of int 
| IOp of int 
| Int of int 
type prog = code list 
type item = 
| Val of code 
| Env of int 
| Code of int
type stack = item list 

match (code, stack,prog) 
(_, _, (Ldi _)::_)
(_, _, (Push _)::_)
    \end{lstlisting} &5&3&11&1&3&30\\
        \hline      
        \newpage
        \begin{lstlisting}[basicstyle=\scriptsize,belowskip=-1.5em,aboveskip=-1em]
(_, _, (Ldi _)::_)
(_, _, (Push _)::_)
(Int(_), _, (IOp _)::_)
\end{lstlisting}
     &20&7&$\sim$1700&3&5&$\sim$11000                    \\ \hline
  \end{longtable}

\end{frame}

\end{comment}

%%%%%%%%%%%%%%%%%%%%%%%%%%%%%%%%%%%%%%%%%%%%%%%%
\begin{frame}[fragile]{Evaluation: matching three boolean values}
\begin{minipage}[c][5cm][t]{0.35\linewidth}
\begin{lstlisting}[basicstyle=\small,belowskip=1em,aboveskip=0em]
(_, false, true) -> 1
(false, true, _) -> 2
(_, _, false) -> 3
(_, _, true) -> 4
\end{lstlisting}
\end{minipage}
\begin{minipage}{0.63\linewidth}
\begin{onlyenv}<1>
Answer of size 6 (in 1.6 seconds)
\begin{lstlisting}[basicstyle=\small,belowskip=0em]
switch $\bullet$[0] with  
| true -> 
    (switch $\bullet$[1] with  
    | true -> 
        (switch $\bullet$[2] with true -> 4 | _ -> 3)
    | _ -> 
        (switch $\bullet$[2] with true -> 1 | _ -> 3))  
| _ -> 
    (switch $\bullet$[1] with  
    | true -> 2   
    | _ -> 
         (switch $\bullet$[2] with true -> 1 | _ -> 3))
\end{lstlisting}
\end{onlyenv}
\begin{onlyenv}<2>
Answer of size 5 (in +0.4 seconds)
\begin{lstlisting}[basicstyle=\small,belowskip=1em]
switch $\bullet$[0] with  
| true -> 
    (switch $\bullet$[2] with  
    | true -> 
       (switch $\bullet$[1] with true -> 4 | _ -> 1)
    | _ -> 3 )  
| _ -> 
    (switch $\bullet$[1] with  
    | true -> 2   
    | _ -> (switch $\bullet$[2] with  true -> 1 | _ -> 3))
\end{lstlisting}
\end{onlyenv}
\begin{onlyenv}<3>
Answer of size 4 (in +0.7 seconds)
\begin{lstlisting}[basicstyle=\small,belowskip=0em]
switch $\bullet$[1] with  
| true -> 
    (switch $\bullet$[0] with  
    | true -> (switch $\bullet$[2] with true -> 4 | _ -> 3)
    | _ -> 2) 
| _ -> 
    (switch $\bullet$[2] with true -> 1 | _ -> 3)
\end{lstlisting}
\end{onlyenv}
\end{minipage}
\end{frame}

%%%%%%%%%%%%%%%%%%%%%%%%%%%%%%%%%%%%%%%%%%%%%%%%
\begin{comment}
\begin{frame}[fragile]{Evaluation (1/3)}
\begin{figure}
\begin{subfigure}[b]{0.48\linewidth}
\begin{lstlisting}[basicstyle=\small,belowskip=1em,aboveskip=0em]
(_, false, true) -> 1
(false, true, _) -> 2
(_, _, false) -> 3
(_, _, true) -> 4
\end{lstlisting}
Answer of size 6 (in 1.6 seconds)
\begin{lstlisting}[basicstyle=\scriptsize,belowskip=0em]
switch $\bullet$[0] with  
| true -> 
    (switch $\bullet$[1] with  
    | true -> 
        (switch $\bullet$[2] with true -> 4 | _ -> 3)
    | _ -> 
        (switch $\bullet$[2] with true -> 1 | _ -> 3))  
| _ -> 
    (switch $\bullet$[1] with  
    | true -> 2   
    | _ -> 
         (switch $\bullet$[2] with true -> 1 | _ -> 3))
\end{lstlisting}
\end{subfigure}
%\hspace{.5cm}
\begin{subfigure}[b]{0.48\linewidth}
Answer of size 5 (in +0.4 seconds)
\begin{lstlisting}[basicstyle=\scriptsize,belowskip=1em]
switch $\bullet$[0] with  
| true -> 
    (switch $\bullet$[2] with  
    | true -> 
       (switch $\bullet$[1] with true -> 4 | _ -> 1)
    | _ -> 3 )  
| _ -> 
  (switch $\bullet$[1] with  
  | true -> 2   
  | _ -> (switch $\bullet$[2] with  true -> 1 | _ -> 3))
\end{lstlisting}
Answer of size 4 (in +0.7 seconds)
\begin{lstlisting}[basicstyle=\scriptsize,belowskip=0em]
switch $\bullet$[1] with  
| true -> 
   (switch $\bullet$[0] with  
   | true -> (switch $\bullet$[2] with true -> 4 | _ -> 3)
   | _ -> 2) 
| _ -> (switch $\bullet$[2] with true -> 1 | _ -> 3)
\end{lstlisting}
\end{subfigure}
\end{figure}
\end{frame}


content
\end{comment}

\begin{frame}[fragile]{Performance problem (1/2)}
\begin{figure}
\begin{subfigure}[b]{0.5\linewidth}
\begin{lstlisting}[basicstyle=\small,language=ocaml]
match a,s,c with
| (_,_,Ldi i::_) -> 1
| (_,_,Push::_)  -> 2
| (Int _,Val (Int _)::_,IOp _::_) -> 3
| (Int _,_,Test (_,_)::c) -> 4
| (Int _,_,Test (_,_)::c) -> 5
| (_,_,Extend::_) -> 6
| (_,_,Search _::_) -> 7
| (_,_,Pushenv::_) -> 8
| (_,Env e::s,Popenv::_) -> 9
| (_,_,Mkclos cc::_) -> 10
| (_,_,Mkclosrec _::_) -> 11
| (Clo (_,_), Val _::_, Apply::_) -> 12
| (_,(Code _::Env _::_),[]) -> 13
| (_,[],[]) -> 14
\end{lstlisting}
\end{subfigure}
\hspace{1cm}
\begin{subfigure}[b]{0.35\linewidth}
Doesn't currently work because the types are too large (11102 examples  generated).\\

%But for a reduced example we can synthesize the answer
\vspace{1cm}
\end{subfigure}
\end{figure}
\end{frame}

\begin{frame}[fragile]{Performance problem (2/2)}
\begin{figure}
\begin{subfigure}[b][6cm][b]{0.45\linewidth}
\begin{lstlisting}[basicstyle=\small]
type code = 
| Push 
| Ldi of int 
| IOp of int 
| Int of int 
type prog = code list 
type item = 
| Val of code 
| Env of int 
| Code of int
type stack = item list 

match (code, stack,prog) with
| (_, _, (Ldi _)::_) -> 1
| (_, _, (Push _)::_) -> 2
\end{lstlisting}
\begin{onlyenv}<1>
\vspace{4mm}
\end{onlyenv}
\begin{onlyenv}<2>
\begin{lstlisting}[basicstyle=\small,aboveskip=-0.5em]
| (Int _, _, (IOp _)::_) -> 3
\end{lstlisting}
\end{onlyenv}
\end{subfigure}
\hspace{.5cm}
\begin{subfigure}[b]{0.45\linewidth}
\begin{overlayarea}{7cm}{6cm}
Example is reduces by 
\begin{itemize}
\item types 
\item clauses
\end{itemize}
For two clauses we require 5 examples \\\vspace{1em}

\begin{onlyenv}<2->
For three clauses and the same types we require 20 examples
\begin{itemize}
\item in 1,5s it gives the  1st answer of size 7
\item in the next half a second it gives the 2nd and the 3rd (last) answers of sizes 6 and 5
\item in the end it in 10s proves that no shorter answers exist
\end{itemize}
%\begin{lstlisting}[basicstyle=\small,aboveskip=-0.5em]
%| (Int _, _, (IOp _)::_) -> 3
%\end{lstlisting}
\end{onlyenv}
\end{overlayarea}
\vspace{1cm}
\end{subfigure}
\end{figure}

\end{frame}

\begin{frame}{Conclusions and Future Work}
Contributions: \contributions
\vspace{1em}

An approach shows correct results on small problems...
 but performs not fast enough\\
\vspace{1em}
 
Areas for possible improvements:
\begin{itemize}
\item Reification for structural constraints can be made faster (lazy computations)
\item Use finite-domain constraints instead of disequality constraints 
\item Memoization is not used (because of presence of disequality constraints)
\item More clever way to generate less examples
\end{itemize}
\vspace{1em}\pause
\begin{center}
{\Huge Thanks!}
\end{center}
\end{frame}

\begin{frame}%[t, allowframebreaks]
\frametitle{References}
\bibliographystyle{amsalpha}
\bibliography{references}
\vspace{1cm}
\end{frame}

\end{document}
